\documentclass[11pt]{article}
\usepackage{amsmath,amsthm,amssymb}
\usepackage[margin=1in]{geometry}
\usepackage{enumitem}
\usepackage{booktabs}
\usepackage{hyperref}

\newtheorem{theorem}{Theorem}
\newtheorem{lemma}[theorem]{Lemma}
\newtheorem{proposition}[theorem]{Proposition}
\newtheorem{corollary}[theorem]{Corollary}
\theoremstyle{definition}
\newtheorem{definition}[theorem]{Definition}
\newtheorem{remark}[theorem]{Remark}
\newtheorem{example}[theorem]{Example}

\title{Closing the structural gap $q \in \{3, 4\}$ for $(4, 3, 1^q)$:\\
       $\dim B = 5$ structurally for all $q \ge 3$}
\author{Clio}
\date{13 May 2026 (very late prove session, 2)}

\begin{document}
\maketitle

\begin{abstract}
The May-13 afternoon paper
\texttt{2026-05-13-4-3-1n-lower-bound-closed.tex} proved
$\dim B^{(\lambda)}_{j_0} V_\lambda = 5 = f^{(3, 2)}$ structurally for
$\lambda = (4, 3, 1^q)$ at $q \ge 5$, with computational verification
at $q \in \{3, 4\}$. The structural binding was the dim input on the
sub-shape $\mu_A = (3, 3, 1^{q-1})$, whose own structural range
required $q - 1 \ge 4$ (dim2-331 Theorem 1.3, propagated from the
May-12 evening dim on $(3, 2, 1^{q-2})$).

The very-late prove session
(\texttt{2026-05-13-very-late-low-q-dim.tex}) closed the
$(3, 2, 1^{q'})$ dim formula structurally for \emph{every} $q' \ge 1$
(Theorem~A) and the $(4, 2, 1^{q'})$ dim formula for every
$q' \ge 2$ (Theorem~B). Substituting Theorem~A as the input to
dim2-331 extends the latter down to $q' \ge 2$ for $(3, 3, 1^{q'})$
(Lemma~\ref{lem:dim2-331-extended} below), unlocking the structural
$\mu_A$ input on $(4, 3, 1^q)$ at $q - 1 \ge 2$.

\medskip
\textbf{Theorem.} For $\lambda = (4, 3, 1^q)$ at every $q \ge 3$,
\[
   \dim B^{(\lambda)}_{j_0} V_\lambda \;=\; 5 \;=\; f^{(3, 2)},
\]
\emph{structurally}. Combined with the May-13 afternoon paper for
$q \ge 5$ and the computational match at $q \in \{3, 4\}$, this gives
the closed-form dim formula $\dim B = f^{\lambda^{(\ge 2)}}$ for the
entire $(4, 3, 1^*)$ family at $q \ge 3$.

\medskip
The proof has \emph{no new technical content}: every ingredient is
already established in the May-13 papers (afternoon-2 dim2-331,
afternoon-3 lower-bound-closed, very-late). The contribution is to
identify that the new $(3, 2, 1^*)$ low-$q$ closure (Theorem~A) is
exactly the input that propagates dim2-331 down to $q' \ge 2$, which
in turn propagates the $(4, 3, 1^*)$ proof down to $q \ge 3$. This is
a single propagation step along the existing SRD ladder.

\medskip
Mod-$P$ verification at $P = 100\,003$, $Q = 11$ confirms every
intermediate dim value.
\end{abstract}

\section{Setup and inputs}\label{sec:setup}

Notation as in the May-13 afternoon papers. $H_q(S_n)$ is the
Iwahori--Hecke algebra over $\mathbb{Q}(q)$ with generators
$T_1, \ldots, T_{n-1}$ obeying $(T_i - q)(T_i + 1) = 0$ and the
standard braid relations. $V_\lambda$ is the seminormal cell module
for $\lambda \vdash n$ in asymmetric Murphy normalisation; $E_j^\pm$
are the $q$- and $(-1)$-eigenspaces of $T_j$. For $1 \le k \le n$,
\[
   R'_k \;:=\; (T_{k-1}{+}1)\cdots(T_1{+}1),
   \qquad
   B^{(\lambda)}_j V_\lambda \;:=\; R'_j\cdots R'_n V_\lambda,
   \qquad
   j_0(\lambda) := \ell(\lambda) + 1.
\]

We cite the following inputs verbatim from the May-13 papers.

\begin{theorem}[\textbf{A} of very-late, \texttt{2026-05-13-very-late-low-q-dim.tex} Theorem 1]\label{thm:A}
For $\nu = (3, 2, 1^p)$ with $p \ge 1$ and $|\nu| = p + 5$,
\[
   \dim B^{(\nu)}_{j_0(\nu)} V_\nu \;=\; 2 \;=\; f^{(2, 1)},
\]
\emph{structurally}.
\end{theorem}

\begin{theorem}[\textbf{B} of very-late, same paper, Theorem 2]\label{thm:B}
For $\nu = (4, 2, 1^p)$ with $p \ge 2$ and $|\nu| = p + 6$,
\[
   \dim B^{(\nu)}_{j_0(\nu)} V_\nu \;=\; 3 \;=\; f^{(3, 1)},
\]
\emph{structurally}.
\end{theorem}

\begin{theorem}[SRD-331 Corollary 4.3 + Lemma 3.2]\label{thm:srd331-reduction}
For $\mu = (3, 3, 1^{q'})$ with $q' \ge 2$, $n_\mu := q' + 6$,
$\ell_\mu := q' + 2$, $j_0(\mu) := \ell_\mu + 1$, and
$\mu_1 := \mu \setminus \{c_1, c_*\} = (3, 2, 1^{q'-1})$
($c_1 = (2, 3)$, $c_* = (q'+2, 1)$, axial distance
$|d| = q' + 2 \ge 4$),
\begin{equation}\label{eq:srd331-reduction}
   B^{(\mu)}_{j_0(\mu)} V_\mu
   \;=\;
   (T_{n_\mu-4}{+}1)(T_{n_\mu-3}{+}1)(T_{n_\mu-2}{+}1)\,
   \Phi_*\!\big(B^{(\mu_1)}_{j_0(\mu_1)} V_{\mu_1}\big),
\end{equation}
where $\Phi_* : V_{\mu_1} \hookrightarrow V_\mu$ is the
$\{c_1, c_*\}$-branching injection: $\Phi_*(v_{T'})$ is the
$T_{n_\mu - 1}$-$(+q)$-eigenvector of the $2$-block at
$\{c_1, c_*\}$. $\Phi_*$ is $H_q(S_{n_\mu - 2})$-equivariant and
$\Phi_*(V_{\mu_1}) \subseteq E_{n_\mu - 1}^+|_{V_\mu}$.
\end{theorem}

\begin{proposition}[Adjacent-injectivity, May-19 + dim2-331 Theorem 3.4]\label{prop:adj-inj}
For any $H_q(S_n)$-module $V$ and any $1 \le a \le b \le n - 2$, the
composition
\[
   (T_a{+}1)(T_{a+1}{+}1)\cdots(T_b{+}1)\big|_{E_{b+1}^+(V)}
   \;:\;
   E_{b+1}^+(V) \;\longrightarrow\; E_a^+(V)
\]
is injective. Each factor $(T_j{+}1) : E_{j+1}^+(V) \to E_j^+(V)$ is
well-defined (Hecke quadratic relation) and injective (May-19
adjacent disjointness $E_j^- \cap E_{j+1}^+ = 0$).
\end{proposition}

\begin{theorem}[Four-branch reduction, May-13-UB Theorem~3.4 + Lemmas 3.5, 3.6]\label{thm:reduction}
For $\lambda = (4, 3, 1^q)$ at every $q \ge 3$,
\begin{equation}\label{eq:four-branch}
   B^{(\lambda)}_{j_0} V_\lambda
   \;=\;
   \mathcal{O}_\lambda\!\left[
     \Phi_A^\lambda\!\big(B^{(\mu_A)}_{j_0(\mu_A)} V_{\mu_A}\big)
     +
     \Phi_B^\lambda\!\big(B^{(\mu_B)}_{j_0(\mu_B)} V_{\mu_B}\big)
   \right],
\end{equation}
where $\mu_A = (3, 3, 1^{q-1})$, $\mu_B = (4, 2, 1^{q-1})$,
$\mathcal{O}_\lambda = (T_{q+2}{+}1)(T_{q+3}{+}1)(T_{q+4}{+}1)(T_{q+5}{+}1)$
(four factors at indices $\ell, \ell+1, n-3, n-2$).
\textbf{The reduction is fully structural at $q \ge 3$:} the third
($\mu_C$-piece) and fourth (same-row piece) summands of the four-branch
decomposition vanish at $q \ge 3$ (May-13-UB Lemma~3.5 + Lemma~3.6),
citing the hook $j$-formulas on $\mu_C = (3, 2, 1^q)$ at $q \ge 1$
(Shift Lemma + hook $j$-formula chain) and on $\nu_2 = (4, 1^{q+1})$
at $q + 1 \ge 2$ (Shift Lemma).
\end{theorem}

\begin{remark}\label{rem:reduction-no-mu-j-needed}
The reduction~\eqref{eq:four-branch} \emph{does not depend} on the
$j$-formulas for the surviving sub-shapes $\mu_A, \mu_B$. It depends
only on the $j$-formulas for the vanishing pieces $\mu_C$ and
$\nu_2$, both of which are hooks and structural for all $q \ge 1$.
This is the key observation enabling the structural upper bound at
$q \ge 3$: the surviving pieces enter only via their \emph{dimensions},
which my new inputs (Lemma~\ref{lem:dim2-331-extended} and
Theorem~\ref{thm:B}) supply at $q \ge 3$.
\end{remark}

\begin{theorem}[Disjoint $\Phi$-supports + image-in-$E^+$, May-13 afternoon-3 Lemmas 2.1, 2.2]\label{thm:disjoint-image}
For $\lambda = (4, 3, 1^q)$ and the branching injections $\Phi_A^\lambda$,
$\Phi_B^\lambda$ at corner pairs $\{c_1, c_3\}, \{c_2, c_3\}$
(where $c_1 = (1, 4)$, $c_2 = (2, 3)$, $c_3 = (q+2, 1)$):
\begin{enumerate}[leftmargin=2em,topsep=0.2em,itemsep=0.1em]
\item $\Phi_A^\lambda(V_{\mu_A}) \cap \Phi_B^\lambda(V_{\mu_B}) = \{0\}$
inside $V_\lambda$ (their seminormal SYT-supports are disjoint
because the unordered cell pairs $\{c_1, c_3\}, \{c_2, c_3\}$
differ).
\item $\Phi_A^\lambda(V_{\mu_A}) \cup \Phi_B^\lambda(V_{\mu_B})
\subseteq E_{n-1}^+|_{V_\lambda}$.
\end{enumerate}
Each $\Phi_X^\lambda$ is itself injective on $V_{\mu_X}$.
\end{theorem}

\section{Extended dim2-331: $\dim B^{(3, 3, 1^{q'})} = 2$ structurally at $q' \ge 2$}\label{sec:ext-dim2-331}

\begin{lemma}[Extended dim2-331]\label{lem:dim2-331-extended}
For $\mu = (3, 3, 1^{q'})$ at every $q' \ge 2$,
\[
   \dim B^{(\mu)}_{j_0(\mu)} V_\mu \;=\; 2,
\]
\emph{structurally}.
\end{lemma}

\begin{proof}
The dim2-331 paper (\texttt{2026-05-13-dim2-331-structural.tex})
proves $\dim B^{(\mu)} = 2$ structurally at $q' \ge 4$ via three
ingredients:
\begin{enumerate}[label=(\roman*),leftmargin=2.2em,topsep=0.2em,itemsep=0.1em]
\item SRD-331 reduction formula~\eqref{eq:srd331-reduction}
  (Theorem~\ref{thm:srd331-reduction}): unconditional at $q' \ge 2$.
\item Non-degeneracy of the corner-pair $2$-block: $|d| = q' + 2 \ge 4$
  at $q' \ge 2$.
\item Dim of the inner space: $\dim B^{(\mu_1)}_{j_0(\mu_1)} V_{\mu_1}
  = 2$ for $\mu_1 = (3, 2, 1^{q'-1})$.
\end{enumerate}
The original dim2-331 used the May-12 evening dim formula on
$\mu_1$, requiring $|\mu_1| \ge 8$, i.e., $q' \ge 4$. Substituting
Theorem~A (very-late, Theorem~\ref{thm:A} above) for the inner dim:
$\dim B^{(\mu_1)} = 2$ for $\mu_1 = (3, 2, 1^{q'-1})$ at
$q' - 1 \ge 1$, i.e., \emph{every} $q' \ge 2$.

With this strictly stronger inner-dim input, the dim2-331 argument
runs unchanged:
\begin{enumerate}[leftmargin=2em,topsep=0.2em,itemsep=0.2em]
\item Upper bound:
$B^{(\mu)} = \mathcal{P}_\mu\,\Phi_*(B^{(\mu_1)})$ (Theorem~\ref{thm:srd331-reduction}),
$\Phi_*$ injective (axial distance $|d| = q' + 2 \ge 4$),
$\dim B^{(\mu_1)} = 2$ (Theorem~A applied at $p = q' - 1 \ge 1$).
Hence
$\dim B^{(\mu)} \le \dim \Phi_*(B^{(\mu_1)}) = 2$.

\item Lower bound: $\Phi_*(B^{(\mu_1)}) \subseteq E_{n_\mu - 1}^+|_{V_\mu}$
is $2$-dimensional. The outer chain
$\mathcal{P}_\mu = (T_{n_\mu-4}{+}1)(T_{n_\mu-3}{+}1)(T_{n_\mu-2}{+}1)$
is the composition of three adjacent factors, indices
$n_\mu - 4, n_\mu - 3, n_\mu - 2$ in increasing order. By
Proposition~\ref{prop:adj-inj} with $a = n_\mu - 4$, $b = n_\mu - 2$,
$\mathcal{P}_\mu$ acts injectively on $E_{n_\mu - 1}^+|_{V_\mu}$.
Hence
$\dim \mathcal{P}_\mu \Phi_*(B^{(\mu_1)}) = \dim \Phi_*(B^{(\mu_1)}) = 2$,
which gives $\dim B^{(\mu)} \ge 2$.
\end{enumerate}
Combining, $\dim B^{(\mu)} = 2$.
\end{proof}

\begin{remark}\label{rem:no-new-ingredient}
Lemma~\ref{lem:dim2-331-extended} introduces no new technical
machinery. The dim2-331 argument was always shape-agnostic in
ingredients (i), (ii), and the adjacent-injectivity chain. The only
shape-specific dependency was the inner-dim input on $\mu_1$.
Improving that input automatically improves the conclusion's range.
\end{remark}

\section{Main theorem: $\dim B^{(4, 3, 1^q)} = 5$ structurally at $q \ge 3$}\label{sec:main}

\begin{theorem}[Main]\label{thm:main}
For $\lambda = (4, 3, 1^q)$ at every $q \ge 3$,
\[
   \dim B^{(\lambda)}_{j_0} V_\lambda \;=\; 5 \;=\; f^{(3, 2)},
\]
\emph{structurally}.
\end{theorem}

\begin{proof}
We assemble the May-13 afternoon-3 proof at the new (lower) threshold.
The structural reduction~\eqref{eq:four-branch} gives
\[
   B^{(\lambda)}_{j_0} V_\lambda
   \;=\;
   \mathcal{O}_\lambda(W),
   \qquad
   W \;:=\;
   \Phi_A^\lambda(B^{(\mu_A)}_{j_0(\mu_A)} V_{\mu_A})
   \;+\;
   \Phi_B^\lambda(B^{(\mu_B)}_{j_0(\mu_B)} V_{\mu_B}).
\]
We will establish $\dim W = 5$ (Step 1) and that $\mathcal{O}_\lambda$
acts injectively on a containing subspace (Steps 2--3); together these
yield $\dim \mathcal{O}_\lambda W = 5$, hence
$\dim B^{(\lambda)}_{j_0} V_\lambda = 5$.

\textbf{Step 1: $\dim W = 5$.}
\begin{itemize}[leftmargin=2em,topsep=0.2em,itemsep=0.1em]
\item By Lemma~\ref{lem:dim2-331-extended} applied at $q' = q - 1 \ge 2$
(i.e., $q \ge 3$):
$\dim B^{(\mu_A)}_{j_0(\mu_A)} V_{\mu_A} = 2$, structurally.
\item By Theorem~\ref{thm:B} applied at $p = q - 1 \ge 2$ (i.e., $q \ge 3$):
$\dim B^{(\mu_B)}_{j_0(\mu_B)} V_{\mu_B} = 3$, structurally.
\item By Theorem~\ref{thm:disjoint-image}(1), the seminormal
$V_\lambda$-supports of $\Phi_A^\lambda(V_{\mu_A})$ and
$\Phi_B^\lambda(V_{\mu_B})$ are disjoint, hence the two images
intersect trivially; their sum is a direct sum. Each $\Phi_X^\lambda$
is itself injective on $V_{\mu_X}$ (distinct seminormal basis vectors
go to vectors with disjoint $2$-element supports).
Hence
\[
   \dim W
   \;=\;
   \dim \Phi_A^\lambda(B^{(\mu_A)}) + \dim \Phi_B^\lambda(B^{(\mu_B)})
   \;=\;
   \dim B^{(\mu_A)} + \dim B^{(\mu_B)}
   \;=\;
   2 + 3 \;=\; 5.
\]
\end{itemize}

\textbf{Step 2: $W \subseteq E_{n-1}^+|_{V_\lambda}$.}
By Theorem~\ref{thm:disjoint-image}(2), both summands of $W$ lie in
$E_{n-1}^+|_{V_\lambda}$, hence so does $W$.

\textbf{Step 3: $\mathcal{O}_\lambda$ is injective on $W$.}
The outer chain is the four-factor product
$\mathcal{O}_\lambda = (T_{q+2}{+}1)(T_{q+3}{+}1)(T_{q+4}{+}1)(T_{q+5}{+}1)$
with indices $a = q + 2 = \ell$ and $b = q + 5 = n - 2$.
Proposition~\ref{prop:adj-inj} applies: $\mathcal{O}_\lambda$ acts
injectively on all of $E_{n-1}^+|_{V_\lambda}$, hence on the
subspace $W$.

\textbf{Conclusion.} By Steps 1 and 3:
\[
   \dim \mathcal{O}_\lambda W \;=\; \dim W \;=\; 5.
\]
By the four-branch reduction~\eqref{eq:four-branch},
$B^{(\lambda)}_{j_0} V_\lambda = \mathcal{O}_\lambda W$. Hence
\[
   \dim B^{(\lambda)}_{j_0} V_\lambda \;=\; \dim \mathcal{O}_\lambda W
   \;=\; 5.
\]
(The argument gives both bounds simultaneously: $\dim B = \dim \mathcal{O}_\lambda W
\le \dim W = 5$ for the upper bound, and $\dim B = \dim \mathcal{O}_\lambda W
= \dim W = 5$ for the lower bound, by injectivity.)
\end{proof}

\begin{remark}[Range of validity, sub-shape audit]\label{rem:audit}
We audit each structural input against the regime $q \ge 3$:
\begin{center}
\renewcommand{\arraystretch}{1.15}
\begin{tabular}{lll}
\toprule
Ingredient & Required & Threshold at $\lambda = (4, 3, 1^q)$ \\
\midrule
Phase~A four-branch decomposition (May-13-UB Thm~3.4) & $q \ge 3$ & $q \ge 3$ \\
$\mu_C$-piece vanishes (May-13-UB Lem~3.5, hook $j$-form on $(3, 2, 1^q)$) & $q \ge 1$ & $q \ge 3$ \\
Same-row vanishes (May-13-UB Lem~3.6, hook $j$-form on $(4, 1^{q+1})$) & $q \ge 1$ & $q \ge 3$ \\
$\dim B^{(\mu_A)}$ structural & $q - 1 \ge 2$ & $q \ge 3$ \\
\hspace{1em} (Lemma~\ref{lem:dim2-331-extended}, via Theorem~A on $\mu_1 = (3,2,1^{q-2})$) & $q - 2 \ge 1$ & $q \ge 3$ \\
$\dim B^{(\mu_B)}$ structural (Theorem~\ref{thm:B}) & $q - 1 \ge 2$ & $q \ge 3$ \\
Disjoint $\Phi$-supports (Theorem~\ref{thm:disjoint-image}(1)) & shape-agnostic & all $q$ \\
$\Phi_X \subseteq E_{n-1}^+$ (Theorem~\ref{thm:disjoint-image}(2)) & $|d_X| \ge 2$ & all $q \ge 1$ \\
Adjacent-injectivity of $\mathcal{O}_\lambda$ (Prop~\ref{prop:adj-inj}) & shape-agnostic & all $q$ \\
\bottomrule
\end{tabular}
\end{center}
Every input is available at $q \ge 3$; the binding constraint is the
$\mu_A$ input via Lemma~\ref{lem:dim2-331-extended}, which itself
binds at $q' = q - 1 \ge 2$ via Theorem~A at $p = q' - 1 = q - 2 \ge 1$.
\end{remark}

\begin{remark}[Comparison to the May-13 afternoon-3 paper]\label{rem:vs-afternoon-3}
The May-13 afternoon-3 paper used the original dim2-331 (binding
$q' \ge 4$) for the $\mu_A$ input, forcing $q - 1 \ge 4$, i.e.,
$q \ge 5$ as the binding threshold. Replacing the inner input with
Theorem~A (which sharpens dim2-331's $\mu_1$-input from $q' - 1 \ge 4$
to $q' - 1 \ge 1$) drops every subsequent threshold by $3$ units:
dim2-331 goes from $q' \ge 4$ to $q' \ge 2$, hence $(4, 3, 1^q)$ goes
from $q \ge 5$ to $q \ge 3$.
\end{remark}

\section{Computational verification}\label{sec:comp}

Mod-$P$ computation at $P = 100\,003$, $Q = 11$ via the
\texttt{compute\_dim\_B.py} script (\texttt{scratch/2026-05-13-multi-corner/}):

\begin{center}
\renewcommand{\arraystretch}{1.15}
\begin{tabular}{llcc}
\toprule
shape & $n$ & $\dim V$ & $\dim B$ \\
\midrule
$(4, 3, 1^3)$ & $10$ & $525$ & $\mathbf{5}$ \\
$(4, 3, 1^4)$ & $11$ & $1100$ & $\mathbf{5}$ \\
$(3, 3, 1^2)$ & $8$ & $56$ & $2$ \\
$(3, 3, 1^3)$ & $9$ & $120$ & $2$ \\
$(4, 2, 1^2)$ & $8$ & $90$ & $3$ \\
$(4, 2, 1^3)$ & $9$ & $189$ & $3$ \\
$(3, 2, 1^1)$ & $6$ & $16$ & $2$ \\
$(3, 2, 1^2)$ & $7$ & $35$ & $2$ \\
\bottomrule
\end{tabular}
\end{center}

The first two rows verify the main theorem at the two new structural
cases $q \in \{3, 4\}$. Rows 3--4 verify the
$\mu_A = (3, 3, 1^{q-1})$ input at $q' \in \{2, 3\}$, matching
Lemma~\ref{lem:dim2-331-extended}. Rows 5--6 verify the
$\mu_B = (4, 2, 1^{q-1})$ input at $p \in \{2, 3\}$, matching
Theorem~\ref{thm:B}. Rows 7--8 verify the $\mu_1 = (3, 2, 1^{q-2})$
input for Lemma~\ref{lem:dim2-331-extended} at $p \in \{1, 2\}$,
matching Theorem~\ref{thm:A}.

\section{Outlook and the SRD-ladder shape audit}\label{sec:outlook}

\begin{remark}[Status of the $r = 2$ same-row club]\label{rem:club-status}
After this paper, the structural $j$-formula club for $r = 2$
length-preserving-corner shapes covers the following stable regimes:
\begin{center}
\renewcommand{\arraystretch}{1.15}
\begin{tabular}{lll}
\toprule
family & structural range & source \\
\midrule
$(3, 2, 1^q)$ & $q \ge 1$ & very-late Theorem A \\
$(4, 2, 1^q)$ & $q \ge 2$ & very-late Theorem B \\
$(3, 3, 1^q)$ & $q \ge 2$ & Lemma~\ref{lem:dim2-331-extended} (this paper) \\
$(4, 3, 1^q)$ & $q \ge 3$ & Theorem~\ref{thm:main} (this paper) \\
$(5, 2, 1^q)$ & $q \ge 5$ & afternoon-3 (only) \\
\bottomrule
\end{tabular}
\end{center}

The next concrete target is $(5, 2, 1^q)$ at $q \in \{3, 4\}$.
Following the same propagation analysis: the SRD-ladder template
(see \texttt{srd-ladder-climb} memory) for $(5, 2, 1^q)$ has
$\mu_A = (4, 2, 1^{q-1})$, $\mu_B = (5, 1^q)$ (hook),
$\mu_C = (4, 1^{q+1})$ (hook), $\nu_1 = (3, 2, 1^q)$.
With very-late Theorem~B closing $\mu_A$ at $q - 1 \ge 2$ and
Theorem~A closing $\nu_1$ at $q \ge 1$, the binding constraint
should drop to $q \ge 3$. Pending: the actual same-row vanishing
and reduction at $q \in \{3, 4\}$ inside the $(5, 2)$-specific
upper-bound paper.
\end{remark}

\begin{remark}[The matching set $\{q : \dim B = f^{\lambda^{(\ge 2)}}\}$]\label{rem:matching}
Combined with the dim2-331-non-monotonicity remark in the
decoration-iso-failed paper, this paper sharpens the picture of the
matching set for $(4, 3, 1^q)$:
\[
   \{q \ge 0 : \dim B^{(4, 3, 1^q)} = f^{(3, 2)} = 5\}
   \;=\; \{3, 4, 5, 6, \ldots\}
   \;\stackrel{?}{=}\; \{q \ge 3\}.
\]
At $q \in \{0, 1, 2\}$, direct computation (decoration-iso-failed
Section~5) shows the failure values $\dim = 4, 7, 12$.
The threshold $q \ge 3$ matches the conjecture
$q_0^{\dim}(\widehat\lambda) = \max(0, |\widehat\lambda| - 2\ell(\widehat\lambda))$
at $\widehat\lambda = (4, 3)$: $|\widehat\lambda| - 2\ell = 7 - 4 = 3$.
\end{remark}

\end{document}
